\documentclass[10pt]{article}
\usepackage[margin=0.9in]{geometry}
\usepackage{amsmath,amssymb,amsthm,mathtools}
\usepackage[hidelinks]{hyperref}
\usepackage{listings}
\lstset{basicstyle=\ttfamily\scriptsize,breaklines=true,columns=fullflexible,
  keepspaces=true,showstringspaces=false}
\newtheorem{theorem}{Theorem}
\newtheorem{lemma}[theorem]{Lemma}
\newcommand{\E}{\mathbb E}
\newcommand{\Var}{\operatorname{Var}}
\title{Strict log-concavity of the normalized Bell function\\
\large A computer-assisted proof of MathDB problem 324014}
\author{Principia Math In Progress}
\date{September 12, 2026}
\begin{document}
\maketitle

\noindent\textbf{Task 1: solve the original problem.}
For real $x>0$, put
\[
 B(x)=e^{-1}\sum_{k=1}^{\infty}\frac{k^x}{k!}.
\]
We prove Conjecture 5.4 of Chen--Guo--Wang
\cite[p.~649]{CGW}. The argument treats every real $x\ge1$.
Its compact-range part uses rigorous interval arithmetic and explicit
infinite-tail bounds; the remaining unbounded range is handled analytically.

\begin{theorem}\label{main}
The function $B(x)^{1/x}$ is strictly log-concave on $[1,\infty)$.
More precisely,
\[
 \left(\frac{\log B(x)}x\right)''<0\qquad(x\ge1).
\]
Consequently the sequence $\{B_n^{1/n}\}_{n\ge1}$ is strictly log-concave.
\end{theorem}

\section{A curvature identity}
Let $L(x)=\log B(x)$ and give the positive integers the probability
distribution
\[
 p_x(k)=\frac{k^x/k!}{\sum_{j\ge1}j^x/j!}.
\]
All differentiated series converge locally uniformly, by the ratio test
after multiplication by any fixed power of $\log k$. Thus
\[
 L'(x)=\E_x\log K,\qquad L''(x)=\Var_x(\log K).
\]
Define
\[
 G(x)=x^2L''(x)-2xL'(x)+2L(x),\qquad H(x)=xL'(x)-L(x).
                                                        \tag{1}
\]
Then
\[
 \left(\frac{L(x)}x\right)''=\frac{G(x)}{x^3},
 \qquad G(x)=x^2\Var_x(\log K)-2H(x).                    \tag{2}
\]
We show $G(x)<-1/4$ for $1\le x\le10^6$, and $G(x)<-x/5$
for $x\ge10^6$.

\section{An elementary estimate on the unbounded range}
Throughout this section $x\ge10^6$. Let $r>1$ be the unique solution of
$r\log r=x$, and set
\[
 W=\log r,\quad s=\frac r{\sqrt{x}},\quad
 F(t)=x\log t-t\log t+t,\quad
 A=\frac{e^{F(r)}}{\sqrt{2\pi r}}.
\]
We will use
\[
 W>11,\qquad s>80.                                      \tag{3}
\]
Indeed $11e^{11}<10^6<12e^{12}$, so $11<W<12$ at $x=10^6$
and $s>1000/12>80$ there. Both $W$ and
$s=e^{W/2}/\sqrt W$ increase subsequently.
Write
\[
 a_k=\frac{k^x/k!}{A},\quad
 M_0=\sum_{k\ge1}a_k,\quad
 M_2=\sum_{k\ge1}a_k\left(\log\frac kr\right)^2.
\]

\begin{lemma}\label{moments}
For $x\ge10^6$,
\[
 2.17s<M_0<3s,\qquad
 M_2<\frac{3.42s+0.01}{x},\qquad
 \E_x\left(\log\frac Kr\right)^2<\frac{1.6}{x}.          \tag{4}
\]
All decimals in this proof denote exact terminating rational numbers.
\end{lemma}

\begin{proof}
Stirling's bounds \cite{Robbins} give, for integers $k\ge1$,
\[
 a_k=\sqrt{\frac rk}\,
       \exp\{F(k)-F(r)-\epsilon_k\},
 \qquad 0<\epsilon_k<\frac1{12k}.                       \tag{5}
\]
Let $I=[19r/20,21r/20]$. Since
$F'(r)=0$ and $-F''(t)=x/t^2+1/t$, on $I$ we have
\[
 \frac{0.9}{s^2}<-F''(t)<\frac{1.21}{s^2}.             \tag{6}
\]
For the upper estimate, use
$400/361+20/209<1.21$ and $W>11$; for the lower,
$400/441>0.9$ suffices.
On $I$, $k>76000$, and hence
\[
 0.97e^{-1.21(k-r)^2/(2s^2)}
 <a_k<
 1.03e^{-0.9(k-r)^2/(2s^2)}.                           \tag{7}
\]
The prefactors follow from
$\sqrt{20/21}(1-0.001)>0.97$ and $\sqrt{20/19}<1.03$.

We use the elementary lattice-sum bound
\[
 \left|\sum_{k\in\mathbb Z}f(k-r)-\int_{\mathbb R}f(t)\,dt\right|
 \le \operatorname{TV}(f).                            \tag{8}
\]
Here $f$ is integrable and of bounded variation, with the lattice sum
absolutely convergent. To prove (8), partition the real line into
intervals $[k-r,k+1-r]$; the discrepancy on an interval is bounded
by its variation. Summing proves the assertion. For a Gaussian the
variation is $2$; for $t^2e^{-\lambda t^2/2}$ it is $8/(e\lambda)$.

Put $R=r/20$. The function
$e^{-1.21t^2/(2s^2)}\mathbf1_{[-R,R]}(t)$ has variation $2$.
Its integral outside $[-R,R]$ before truncation is at most
\[
 \frac{40s}{1.21\sqrt{x}}e^{-1.21x/800}<0.001s.         \tag{9}
\]
This follows, for example, from $\sqrt{x}\ge1000$,
$1.21x/800>4$, and $(40/1210)e^{-4}<0.001$.
All lattice points in this truncated interval have $k\ge1$.
Equations (7)--(9), $\sqrt{2\pi}>2.5$, and $s>80$ yield
\[
 \sum_{k\in I}a_k
 >0.97\left(\frac{2.5}{1.1}-0.001-\frac2{80}\right)s
 >2.17s.                                             \tag{10}
\]
Similarly, $\sqrt{2\pi}<2.51$ and $\sqrt{0.9}>0.94$ give
\[
 \sum_{k\in I}a_k
 <1.03\left(\frac{2.51}{0.94}+\frac2{80}\right)s
 <2.8s.                                              \tag{11}
\]

On $I$, the mean value theorem gives
$|\log(k/r)|\le20|k-r|/(19r)$. Apply (8) to
$t^2e^{-0.9t^2/(2s^2)}$. Since $0.9^{3/2}>0.85$ and $e>2.7$,
\begin{align*}
 \sum_{k\in I}a_k\left(\log\frac kr\right)^2
 &<\frac{1.03(20/19)^2}{r^2}
       \left(\frac{\sqrt{2\pi}}{0.9^{3/2}}s^3+
                          \frac8{0.9e}s^2\right)\\
 &<\frac{1.03(20/19)^2s^3}{r^2}
       \left(\frac{2.51}{0.85}+
                          \frac8{0.9\cdot2.7\cdot80}\right)
 <\frac{3.42s}{x}.                                    \tag{12}
\end{align*}

We now bound both omitted tails explicitly. For $1\le k<19r/20$,
$-F''(t)\ge x/r^2$ on $[k,r]$, so
\[
 F(k)-F(r)\le-\frac{x}{800}<-\frac{x}{1000}.
\]
Equation (5), the number of such integers being less than $r$, gives
left-tail mass at most $r^{3/2}e^{-x/1000}$ and left-tail second
logarithmic moment at most $r^{3/2}W^2e^{-x/1000}$.
As $r<x$ and $(\log x)^2<x$ for $x\ge10^6$, these are bounded by
$x^{3/2}e^{-x/1000}$ and $x^{5/2}e^{-x/1000}$, respectively.

At $T=21r/20$, (6) gives
\[
 F(T)-F(r)\le-\frac{9x}{8000}<-\frac{x}{1000},
 \qquad F'(T)\le-\frac{9W}{200}.
\]
Concavity therefore bounds the right-tail weights by
$e^{-x/1000}e^{-(9W/200)(k-T)}$. Enumerate the integers $k>T$
as $k=T+z+m$, where $0<z\le1$ and $m=0,1,\ldots$.
Since $e^{-9W/200}<2/3$ and
$0<\log(k/r)\le(k-r)/r\le m+2$, the right-tail mass is at most
$3e^{-x/1000}$ and its second logarithmic moment is at most
\[
 e^{-x/1000}\sum_{m\ge0}(m+2)^2(2/3)^m
       =66e^{-x/1000}.
\]
Thus the combined tail mass is at most $2x^2e^{-x/1000}$ and its
second moment is at most $2x^3e^{-x/1000}$. Both are less than
$0.01/x$: indeed
\[
 2x^4e^{-x/1000}<0.01\qquad(x\ge10^6).                 \tag{13}
\]
The left side decreases on this range; at $x=10^6$ the inequality
follows from $\log200<6$, $\log10^6<14$, and $6+4\cdot14<1000$.
Combining these estimates with (10)--(12) proves the first two
claims of (4). Finally
\[
 \frac{M_2}{M_0}
 <\frac{3.42+0.01/80}{2.17x}<\frac{1.6}{x},
\]
which proves the last claim.
\end{proof}

\begin{lemma}\label{large}
For $x\ge10^6$, $G(x)<-x/5$.
\end{lemma}
\begin{proof}
By (4) and $3/\sqrt{2\pi W}<1$,
\[
 B(x)=e^{-1}A M_0<e^{F(r)-1},\qquad L(x)<F(r)-1.
\]
Cauchy--Schwarz and (4) give
$L'(x)>\log r-\sqrt{1.6/x}$, and hence
\[
 H(x)>x-r-\sqrt{1.6x}+1
 >\left(1-\frac1{11}-\frac1{500}\right)x+1>0.9x.
\]
We used $x\ge10^6>1.6\cdot500^2$.
Also $\Var_x(\log K)\le\E_x(\log(K/r))^2<1.6/x$.
Equation (2) now gives $G(x)<1.6x-1.8x=-x/5$.
\end{proof}

\section{A rigorous certificate for the compact range}
\begin{lemma}\label{compact}
For every real $x\in[1,10^6]$, $G(x)<-1/4$.
\end{lemma}

The proof of this lemma is the finite interval computation below.
We give its formulas, all truncation bounds, the exact cover, and
the complete executable code in the appendix.

For any positive integer $k_0$, let
\[
 t_k=\log(k/k_0),\quad
 v_k=\exp\{xt_k+\log(k_0!)-\log(k!)\},\quad
 S_j=\sum_{k\ge1}v_kt_k^j\quad(j=0,1,2).
\]
Direct substitution in (1) gives the cancellation-resistant identity
\[
 G=x^2\left(\frac{S_2}{S_0}-\left(\frac{S_1}{S_0}\right)^2\right)
 -2x\frac{S_1}{S_0}-2\log(k_0!)+2\log S_0-2.           \tag{14}
\]
For an exact rational interval $[a,b]$, the code chooses integers
$1\le\ell\le k_0\le h$ and encloses the sums over $[\ell,h]$
using Arb real balls. Window selection uses floating point but is
not relied upon for an inequality: the following two ratios are
rigorously checked to be less than $1$.

If $\ell>1$, let
\[
 \rho_L=\ell(1-1/\ell)^a.
\]
The ratio $v_{k-1}/v_k=k(1-1/k)^x$ increases with $k$ and
decreases with $x$. Thus the left omitted mass is at most
\[
 T_L=v_\ell^{+}\frac{\rho_L}{1-\rho_L}.                 \tag{15}
\]
Here a superscript $+$ means an outward upper bound valid on $[a,b]$.
Its contributions to $S_0,S_1,S_2$ lie, respectively, in
$[0,T_L]$, $[-T_L\log k_0,0]$, and $[0,T_L(\log k_0)^2]$.
For $\ell=1$ the left tail is empty.

For the right tail put
\[
 \rho_R=\frac{(1+1/h)^b}{h+1},\quad c=\log(h/k_0),\quad
 R_0=\frac{\rho_R}{1-\rho_R},\quad
 R_1=\frac{\rho_R}{(1-\rho_R)^2},\quad
 R_2=\frac{\rho_R(1+\rho_R)}{(1-\rho_R)^3}.
\]
The adjacent ratio $v_{k+1}/v_k$ decreases with $k$ and increases
with $x$. Furthermore $\log((h+m)/k_0)\le c+m/h$ for $m\ge1$.
The right contributions therefore lie in the nonnegative intervals
with upper endpoints
\[
 v_h^+R_0,\qquad v_h^+(cR_0+R_1/h),\qquad
 v_h^+(c^2R_0+2cR_1/h+R_2/h^2).                        \tag{16}
\]
These enclose the entire infinite sums, including all logarithmic moments.

The initial closed intervals have endpoints
$i^2/64$, $i=8,\ldots,8000$. An interval is accepted only if
Arb proves that the upper bound from (14)--(16) is strictly less
than $-1/4$; otherwise it is bisected at its exact rational midpoint.
At 128-bit precision the run accepts \textbf{8001 intervals},
using \textbf{9 bisections}. It then checks, in exact rational
arithmetic, that the intervals form a contiguous cover from $1$ to $10^6$.
The complete record is \texttt{curvature-certificate-million.json}.
This proves Lemma~\ref{compact}.

\begin{proof}[Proof of Theorem~\ref{main}]
Lemmas~\ref{large} and \ref{compact} cover $[1,\infty)$.
Equation (2) proves the stated strict differential inequality.
Strict concavity, restricted to successive positive integers, gives
$2\log(B_n^{1/n})>\log(B_{n-1}^{1/(n-1)})
 +\log(B_{n+1}^{1/(n+1)})$ for $n\ge2$.
\end{proof}

\section{Verification and scope}
The compact-range proof was run with Python 3.13.14 and
\texttt{python-flint} 0.9.0. Reproduce it with
\begin{lstlisting}
python certify_curvature.py --end 1000000 --prec 128 \
    --output curvature-certificate-million.json
python check_analytic_constants.py
\end{lstlisting}
The second command verifies the 18 elementary constant comparisons
used above; their formulas are already included in the proof.
The first command is an essential part of the proof, not a sample-point
experiment. It relies on the enclosure guarantees of Arb, accessed
through \texttt{python-flint} \cite{Arb,pythonflint}.
The source code and both JSON records accompany this document.

Current eligibility checks, including the misleading MathDB reference
annotation, are documented separately in \texttt{literature-audit.md}.
This is a locally checked computer-assisted proof; no external referee
assessment is claimed.

\clearpage
\section{Follow-up tasks}
\noindent\textbf{Task 2: determine the full positive domain.}
Extend strict log-concavity of $B(x)^{1/x}$ to $0<x<1$, or determine
the maximal interval of positive real arguments on which it holds.

\medskip
\noindent\textbf{Task 3: obtain a proof using analysis alone.}
Replace the compact-range certificate by a direct inequality between
$\Var_x(\log K)$ and $xL'(x)-L(x)$ that is valid for every $x\ge1$.

\medskip
\noindent\textbf{Task 4: classify the Poisson parameter.}
For
$B_\lambda(x)=e^{-\lambda}\sum_{k\ge1}k^x\lambda^k/k!$,
determine all $\lambda>0$ for which $B_\lambda(x)^{1/x}$ is
strictly log-concave for every $x\ge1$.

\medskip
\noindent\textbf{Task 5: obtain sharp curvature estimates.}
Derive a uniform expansion of $G(x)$ in terms of the positive solution
$W e^W=x$, with an explicit remainder, and determine optimal constants
in two-sided bounds for $-(L(x)/x)''$.

\begin{thebibliography}{9}
\bibitem{CGW}
W. Y. C. Chen, J. J. F. Guo and L. X. W. Wang,
Zeta functions and the log-behavior of combinatorial sequences,
\emph{Proc. Edinburgh Math. Soc.} 58 (2015), 637--651.
\href{https://doi.org/10.1017/S0013091515000036}{doi:10.1017/S0013091515000036}.
Conjecture 5.4, p.~649; author manuscript
\href{https://arxiv.org/abs/1208.5213}{arXiv:1208.5213}.
\bibitem{Robbins}
H. Robbins, A remark on Stirling's formula,
\emph{Amer. Math. Monthly} 62 (1955), 26--29.
\href{https://doi.org/10.2307/2308012}{doi:10.2307/2308012}.
\bibitem{Arb}
F. Johansson, Arb: efficient arbitrary-precision midpoint-radius
interval arithmetic, \emph{IEEE Trans. Computers} 66 (2017), 1281--1292.
\href{https://arxiv.org/abs/1611.02831}{arXiv:1611.02831}.
\bibitem{pythonflint}
The python-flint contributors,
\emph{python-flint documentation: real balls},
\url{https://python-flint.readthedocs.io/en/latest/arb.html}.
\end{thebibliography}

\clearpage
\appendix
\section{Complete compact-range certifier}
\lstinputlisting[language=Python]{certify_curvature.py}
\end{document}
